\section{Cronograma}

El cronograma del proyecto se estructura en 6 fases principales, alineadas con la metodología Design Science Research (DSR) y el enfoque ágil Scrum descrito en la sección de metodología. El proyecto tiene una duración estimada de 6 meses (Diciembre 2024 - Mayo 2025), distribuidos en sprints de 2 semanas y fases de validación.

\subsection{Diagrama de Gantt}

\begin{table}[H]
\centering
\setlength{\tabcolsep}{3pt}
\scriptsize
\begin{tabular}{|l|c|c|c|c|c|c|}
\hline
\rowcolor{gray!30}
\textbf{Cronograma de la tesis} & \multicolumn{6}{c|}{\textbf{2024-2025}} \\
\hline
\rowcolor{gray!30}
\textbf{Actividades / Fases} & \textbf{D} & \textbf{E} & \textbf{F} & \textbf{M} & \textbf{A} & \textbf{M} \\
\hline
\multicolumn{7}{|l|}{\cellcolor{gray!20}\textbf{Fase 1: Planeación y Revisión Teórica}} \\
\hline
1. Planteamiento del problema y objetivos de la investigación & \cellcolor{gray!50} & \cellcolor{gray!50} & & & & \\
\hline
2. Marco teórico & \cellcolor{gray!50} & \cellcolor{gray!50} & & & & \\
\hline
3. Revisión de la literatura & \cellcolor{gray!50} & \cellcolor{gray!50} & & & & \\
\hline
4. Formulación de hipótesis y justificación & \cellcolor{gray!50} & & & & & \\
\hline
\multicolumn{7}{|l|}{\cellcolor{gray!20}\textbf{Fase 2: Análisis de Requerimientos (Sprint 0)}} \\
\hline
5. Definición de KPIs y requisitos de hardware & & \cellcolor{gray!50} & & & & \\
\hline
6. Análisis de integración SDK volante y API IA & & \cellcolor{gray!50} & & & & \\
\hline
7. Creación de matriz de trazabilidad inicial & & \cellcolor{gray!50} & & & & \\
\hline
\multicolumn{7}{|l|}{\cellcolor{gray!20}\textbf{Fase 3: Diseño de Arquitectura y UI/UX (Sprint 1)}} \\
\hline
8. Diseño de diagramas de arquitectura del sistema & & & \cellcolor{gray!50} & & & \\
\hline
9. Diseño de interfaz gráfica y mockups & & & \cellcolor{gray!50} & & & \\
\hline
10. Elaboración de guion de pruebas de usuario & & & \cellcolor{gray!50} & & & \\
\hline
\multicolumn{7}{|l|}{\cellcolor{gray!20}\textbf{Fase 4: Desarrollo del Prototipo (Sprints 2-4)}} \\
\hline
11. Sprint 2: Programación en Unity + integración volante & & & \cellcolor{gray!50} & \cellcolor{gray!50} & & \\
\hline
12. Sprint 3: Implementación módulo IA (ML+PLN) & & & & \cellcolor{gray!50} & & \\
\hline
13. Sprint 4: Optimización y generación de ejecutable & & & & \cellcolor{gray!50} & \cellcolor{gray!50} & \\
\hline
14. Ajustes y correcciones iterativas & & & & & \cellcolor{gray!50} & \\
\hline
\multicolumn{7}{|l|}{\cellcolor{gray!20}\textbf{Fase 5: Validación y Pruebas (Sprint 5)}} \\
\hline
15. Validación técnica (KPIs de rendimiento) & & & & & \cellcolor{gray!50} & \\
\hline
16. Ejecución de matriz de trazabilidad & & & & & \cellcolor{gray!50} & \\
\hline
17. Aplicación de pruebas de usuario y cuestionario SUS & & & & & \cellcolor{gray!50} & \\
\hline
18. Análisis de datos y resultados & & & & & & \cellcolor{gray!50} \\
\hline
\multicolumn{7}{|l|}{\cellcolor{gray!20}\textbf{Fase 6: Cierre y Entrega}} \\
\hline
19. Redacción de conclusiones e implicaciones & & & & & & \cellcolor{gray!50} \\
\hline
20. Revisión y corrección del documento final & & & & & & \cellcolor{gray!50} \\
\hline
21. Entrega del informe final de tesis & & & & & & \cellcolor{gray!50} \\
\hline
\end{tabular}
\caption{Cronograma de actividades del proyecto de tesis (Diciembre 2024 - Mayo 2025)}
\label{tab:cronograma}
\end{table}

\vspace{2mm}
\scriptsize{\textit{Nota: D=Diciembre, E=Enero, F=Febrero, M=Marzo, A=Abril, M=Mayo}}

\subsection{Descripción de Fases y Productos Esperados}

\textbf{Fase 1: Planeación y Revisión Teórica (Diciembre 2024 - Enero 2025)}
\begin{itemize}
    \item \textbf{Duración:} 6 semanas
    \item \textbf{Responsable:} Tesista
    \item \textbf{Productos esperados:} Marco teórico completo, planteamiento del problema validado, objetivos e hipótesis definidos
\end{itemize}

\textbf{Fase 2: Análisis de Requerimientos - Sprint 0 (Enero 2025)}
\begin{itemize}
    \item \textbf{Duración:} 2 semanas
    \item \textbf{Responsable:} Tesista con asesoría técnica
    \item \textbf{Productos esperados:} Documento de especificación de requisitos, tabla de KPIs meta (Latencia $<$ 200ms, FPS $>$ 30), matriz de trazabilidad inicial
\end{itemize}

\textbf{Fase 3: Diseño de Arquitectura y UI/UX - Sprint 1 (Febrero 2025)}
\begin{itemize}
    \item \textbf{Duración:} 2 semanas
    \item \textbf{Responsable:} Tesista
    \item \textbf{Productos esperados:} Diagramas de arquitectura del sistema (Unity-Hardware-Nube), mockups de interfaz, guion de prueba de usuario
\end{itemize}

\textbf{Fase 4: Desarrollo del Prototipo - Sprints 2-4 (Febrero - Abril 2025)}
\begin{itemize}
    \item \textbf{Duración:} 10 semanas (3 sprints de 2 semanas + 4 semanas de ajustes)
    \item \textbf{Responsable:} Tesista
    \item \textbf{Productos esperados:} 
    \begin{itemize}
        \item Sprint 2: Módulo de Unity funcional con integración del volante háptico
        \item Sprint 3: Módulo de IA implementado (SDK de Gemini/Firebase, ML y PLN)
        \item Sprint 4: Archivo ejecutable (.exe) del prototipo funcional optimizado
    \end{itemize}
\end{itemize}

\textbf{Fase 5: Validación y Pruebas - Sprint 5 (Abril - Mayo 2025)}
\begin{itemize}
    \item \textbf{Duración:} 4 semanas
    \item \textbf{Responsable:} Tesista con participación de adultos mayores
    \item \textbf{Productos esperados:} Informe de rendimiento técnico (KPIs), matriz de trazabilidad completada, reporte de usabilidad con resultados del cuestionario SUS, análisis de datos cuantitativos y cualitativos
\end{itemize}

\textbf{Fase 6: Cierre y Entrega (Mayo 2025)}
\begin{itemize}
    \item \textbf{Duración:} 2 semanas
    \item \textbf{Responsable:} Tesista con revisión del asesor
    \item \textbf{Productos esperados:} Documento final de tesis completo, presentación de resultados
\end{itemize}

\subsection{Notas Importantes}

\begin{itemize}
    \item El cronograma contempla tiempos para ajustes y correcciones iterativas entre las fases 4 y 5, coherente con la metodología ágil Scrum.
    \item La validación técnica y de usuario (Fase 5) se realizará con el prototipo ejecutable ya compilado.
    \item Se ha reservado tiempo específico para la coordinación con adultos mayores de la Sierra Sur de Oaxaca durante la fase de pruebas.
    \item La duración total del proyecto (6 meses) es realista considerando la complejidad de la integración de hardware, software y modelos de IA.
\end{itemize}
